\section*{Chapitre \ref{ch:g11:lithosphere-asthenosphere} --- Lithosphère et asthénosphère}

\begin{solution}{exo:g11:lithosphere-asthenosphere:1}
La \omterm{def:g11:lithosphere-asthenosphere:lithosphere}{lithosphère} est l'enveloppe externe rigide de la Terre, formée de la
croûte et du manteau le plus superficiel~; l'\omterm{def:g11:lithosphere-asthenosphere:lithosphere}{asthénosphère} est le manteau
plus chaud et moins résistant situé dessous, qui s'écoule lentement sur de
longues durées. La croûte océanique, la croûte continentale et le manteau le
plus superficiel appartiennent à la \omterm{def:g11:lithosphere-asthenosphere:lithosphere}{lithosphère}~; le noyau externe n'en fait
pas partie.
\end{solution}

\begin{solution}{exo:g11:lithosphere-asthenosphere:2}
L'iceberg s'enfonce jusqu'à déplacer son propre poids d'eau~: la fraction
immergée vaut $910/1025 \approx 0.888$, environ \qty{89}{\percent}. La
fraction émergée vaut environ \qty{11.2}{\percent}~; la hauteur totale est
donc de $20/0.112 \approx \qty{180}{m}$, dont environ \qty{160}{m} sous
l'eau.
\end{solution}

\begin{solution}{exo:g11:lithosphere-asthenosphere:3}
$15 + 26 \times 3 = \qty{93}{\celsius}$. À \qty{100}{km}, le même gradient
donnerait plus de \qty{2600}{\celsius}, bien au-dessus de la température de
fusion de la roche du manteau~; or les \omterm{def:g11:seismic-structure:body-waves}{ondes S} traversent le manteau
supérieur, qui est donc solide. Le gradient doit devenir bien plus faible en
profondeur.
\end{solution}

\begin{solution}{exo:g11:lithosphere-asthenosphere:4}
La tranche de 0 à \qty{500}{m} au-dessus du niveau de la mer, avec environ
\qty{16}{\percent} de la surface, et la tranche de 4000 à \qty{4500}{m}
sous le niveau de la mer, avec environ \qty{11}{\percent} (la tranche de
4500 à \qty{5000}{m} sous le niveau de la mer est presque aussi étendue).
\end{solution}

\begin{solution}{exo:g11:lithosphere-asthenosphere:5}
Pour $P = 0$~: \qty{1121}{\celsius}. À \qty{3}{GPa}~: $-5.104 \times 9 +
132.899 \times 3 + 1120.661 \approx -45.9 + 398.7 + 1120.7 \approx
\qty{1473}{\celsius}$.
\end{solution}

\begin{solution}{exo:g11:lithosphere-asthenosphere:6}
Environ \qty{730}{\celsius} à \qty{50}{km} (\omterm{def:g11:lithosphere-asthenosphere:lithosphere}{lithosphère}), environ
\qty{1370}{\celsius} à \qty{150}{km} et \qty{1410}{\celsius} à
\qty{250}{km} (\omterm{def:g11:lithosphere-asthenosphere:lithosphere}{asthénosphère}). La température augmente le plus entre 50 et
\qty{150}{km}, de plus de \qty{600}{\celsius}, parce que la plus grande
partie de cet intervalle se trouve dans la \omterm{def:g11:lithosphere-asthenosphere:lithosphere}{lithosphère}~; entre 150 et
\qty{250}{km}, elle n'augmente que d'environ \qty{45}{\celsius}.
\end{solution}

\begin{solution}{exo:g11:lithosphere-asthenosphere:7}
Racine~: $3 \times 2830/550 \approx \qty{15.4}{km}$. Surplus d'épaisseur~:
$3 \times 3380/550 \approx \qty{18.4}{km}$, c'est-à-dire les \qty{3}{km} de
la chaîne plus sa racine.
\end{solution}

\begin{solution}{exo:g11:lithosphere-asthenosphere:8}
La croûte est plus épaisse de $55 - 41.1 = \qty{13.9}{km}$, et $\Delta H = h
\times 3380/550$, donc $h = 13.9 \times 550/3380 \approx \qty{2.3}{km}$.
\end{solution}

\begin{solution}{exo:g11:lithosphere-asthenosphere:9}
Diminution~: $(4.470 - 4.419)/4.470 \approx \qty{1.1}{\percent}$. Saut~:
$(4.644 - 4.419)/4.419 \approx \qty{5.1}{\percent}$. Avec la profondeur, la
roche est de plus en plus comprimée, ce qui accélère les ondes sismiques~;
une diminution signifie qu'autre chose l'emporte~: la roche est chaude,
proche de la fusion, et moins résistante.
\end{solution}

\begin{solution}{exo:g11:lithosphere-asthenosphere:10}
(1) $(1345 - 610)/60 \approx \qty{12}{\celsius}$ par kilomètre. (2) C'est
environ la moitié du gradient en surface. Une partie de la chaleur est
produite à l'intérieur de la croûte supérieure~; cette chaleur doit elle
aussi traverser les couches supérieures pour s'échapper, ce qui rend le
gradient plus fort près de la surface. Sous la croûte, seule la chaleur
venue de plus bas traverse la roche, et le gradient est plus faible.
\end{solution}

\begin{solution}{exo:g11:lithosphere-asthenosphere:11}
À \qty{200}{km}~: $1300 + 90 = \qty{1390}{\celsius}$~; à \qty{400}{km}~:
$1300 + 180 = \qty{1480}{\celsius}$. L'augmentation est de
\qty{90}{\celsius} sur \qty{200}{km}. À travers la \omterm{def:g11:lithosphere-asthenosphere:lithosphere}{lithosphère}, la
température augmente de
$1345 - 15 = \qty{1330}{\celsius}$ sur \qty{100}{km} seulement~: environ
trente fois plus par kilomètre.
\end{solution}

\begin{solution}{exo:g11:lithosphere-asthenosphere:12}
(1) $118 + 12.3\,z = 1300 + 0.45\,z$ donne $z = 1182/11.85 \approx
\qty{100}{km}$. (2) $118 + 21.5\,z = 1300 + 0.45\,z$ donne $z = 1182/21.05
\approx \qty{56}{km}$. (3) L'ouest, plus chaud, a une \omterm{def:g11:lithosphere-asthenosphere:lithosphere}{lithosphère} bien plus
mince, comme le montrent les mesures (environ \qty{40}{km} là où elle est la
plus mince)~; la \omterm{def:g11:lithosphere-asthenosphere:lithosphere}{lithosphère} très épaisse du milieu du continent demande un
gradient encore plus faible.
\end{solution}

\begin{solution}{exo:g11:lithosphere-asthenosphere:13}
(1) D'après la proposition, $\Delta H = h\,\rho_m/(\rho_m - \rho_c)$, donc $h =
\Delta H\,(\rho_m - \rho_c)/\rho_m$. (2) $550/3380 \approx 0.16$~: environ un
sixième de tout surplus de croûte apparaît sous forme de hauteur, cinq
sixièmes sous forme de racine. (3) Pour l'iceberg, la fraction émergée vaut
$(1025 - 910)/1025 \approx 0.11$~: c'est la même règle, (masse volumique du
fluide $-$ masse volumique du corps qui flotte)/masse volumique du fluide.
\end{solution}

\begin{solution}{exo:g11:lithosphere-asthenosphere:14}
(1) $1489 - 1345 = \qty{144}{\celsius}$. (2) Un peu de liquide se rassemble
en films minces entre les grains de la roche. Les grains ne sont alors plus
solidement imbriqués~: la roche résiste moins au cisaillement, si bien que
les \omterm{def:g11:seismic-structure:body-waves}{ondes S} ralentissent, et elle peut se déformer et s'écouler bien plus
facilement.
\end{solution}

\begin{solution}{exo:g11:lithosphere-asthenosphere:15}
(1) Non. Araser les terres amincirait un peu la croûte continentale, mais
elle resterait épaisse et légère, et flotterait toujours bien plus haut que
la croûte mince et dense des fonds océaniques~; les océans rempliraient
toujours les bassins situés entre les continents. (2) Il existe deux sortes
de croûte, d'épaisseurs et de masses volumiques différentes~: une croûte
continentale épaisse et légère, et une croûte océanique mince et dense.
\end{solution}

\begin{solution}{pb:g11:lithosphere-asthenosphere:1}
\textbf{1.} $5.01 - 0.67 = \qty{4.34}{km}$.

\textbf{2.} $(3380 - 2830)/3380 \approx \qty{16}{\percent}$.

\textbf{3.} La fraction $2830/3380 \approx 0.84$ se trouverait sous la
surface~; l'iceberg s'enfonce un peu plus, $910/1025 \approx 0.89$, parce que
la masse volumique de la glace est plus proche de celle de l'eau que celle
de la croûte ne l'est de celle du manteau.

\textbf{4.} Pendant les secondes du passage d'une onde sismique, ou la durée
d'un séisme, le manteau est un solide~; ce n'est que sur de très longues
durées que l'\omterm{def:g11:lithosphere-asthenosphere:lithosphere}{asthénosphère} s'écoule et laisse flotter la croûte.

\textbf{5.} Sous le \omterm{def:g11:seismic-structure:moho}{Moho} le plus profond, les deux colonnes sont faites de
la même roche du manteau aux mêmes profondeurs, et elles y pèsent autant~;
toutes les différences se trouvent au-dessus de cette profondeur et doivent
être prises en compte.

\textbf{6.} Le continent moyen~: de la croûte depuis \qty{0.67}{km}
au-dessus jusqu'à \qty{40.4}{km} sous le niveau de la mer, puis du manteau
jusqu'à \qty{70}{km}. Le plateau~: de la croûte depuis \qty{5.01}{km}
au-dessus du niveau de la mer jusqu'à un \omterm{def:g11:seismic-structure:moho}{Moho} plus profond de la valeur de
la racine, puis du manteau jusqu'à \qty{70}{km}.

\textbf{7.} La colonne du plateau contient $h + r$ de croûte en plus, la
colonne du continent contient à la place du manteau sur l'épaisseur $r$~:
$\rho_c\,(h + r) = \rho_m\,r$,
donc $r = h\,\rho_c/(\rho_m - \rho_c)$ et $\Delta H = h + r =
h\,\rho_m/(\rho_m - \rho_c)$.

\textbf{8.} $3380/550 \approx 6.15$.

\textbf{9.} $\Delta H = 4.34 \times 6.15 \approx \qty{26.7}{km}$~; la racine
vaut $26.7 - 4.34 \approx \qty{22.3}{km}$.

\textbf{10.} Le \omterm{def:g11:seismic-structure:moho}{Moho} du continent moyen se trouve à $41.1 - 0.67 \approx
\qty{40.4}{km}$ sous le niveau de la mer~; celui du plateau à $40.4 + 22.3 \approx
\qty{62.8}{km}$ sous le niveau de la mer.

\textbf{11.} $41.1 + 26.7 \approx \qty{67.8}{km}$.

\textbf{12.} Un kilomètre de roche au-dessus d'un mètre carré a une masse de
$\rho \times \qty{1000}{m^3}$~: \num{2830} tonnes de croûte ou \num{3380}
tonnes de manteau. Continent~: $41.1 \times 2830 + (70 - 40.43) \times 3380
\approx \num{116300} + \num{99900} \approx \num{216300}$ tonnes. Plateau~:
$67.77 \times 2830 + (70 - 62.76) \times 3380 \approx \num{191800} +
\num{24500} \approx \num{216300}$ tonnes. Les deux sont égales.

\textbf{13.} $41.1 - 0.67 \times 6.15 \approx \qty{37.0}{km}$.

\textbf{14.} $67.8/41.1 \approx 1.65$ fois.

\textbf{15.} $3300/470 \approx 7.02$, $\Delta H \approx 4.34 \times 7.02
\approx \qty{30.5}{km}$, et une croûte d'environ \qty{71.6}{km}. Une
variation d'environ \qty{2}{\percent} de la masse volumique du manteau
modifie l'épaisseur d'environ \qty{4}{km}~: le résultat dépend nettement des
masses volumiques, mais reste proche de \qty{70}{km}.

\textbf{16.} Le modèle donne environ \qty{68}{km}, les ondes sismiques
environ \qty{70}{km}~: ils concordent à quelques pour cent près.

\textbf{17.} $h = 3.23 - 0.67 = \qty{2.56}{km}$, $\Delta H = 2.56 \times
6.15 \approx \qty{15.7}{km}$, et une croûte d'environ \qty{56.8}{km},
quelque \qty{7}{km} de plus que la mesure.

\textbf{18.} Plus mince. Une colonne dont la partie supérieure est plus
légère que la moyenne pèse moins pour la même épaisseur~: elle flotte donc
plus haut. Une croûte plus légère a besoin d'une racine plus petite pour
atteindre une hauteur donnée. La vraie croûte peut donc être plus mince que
celle du modèle, comme le montre la mesure.

\textbf{19.} Sur de longues durées, l'\omterm{def:g11:lithosphere-asthenosphere:lithosphere}{asthénosphère} s'écoule~: elle s'écarte
sous une colonne lourde et revient sous une colonne légère, jusqu'à ce que
les colonnes s'équilibrent.

\textbf{20.} \textbf{La croûte sous le Tibet est épaisse d'environ
\qty{70}{km}}, bien plus que les \qty{41}{km} d'un continent moyen~: le
plateau culmine à cinq kilomètres parce que sa croûte épaisse et légère
flotte sur le manteau plus dense, avec une racine de plus de \qty{20}{km}.
\end{solution}
